\documentclass[10pt,a4paper]{amsart}
\usepackage{amsmath,amssymb,amsthm,graphicx}
% Minimal class-compatibility preamble for a reader excerpt of
%   Bhimarthi Ravinder, "Demazure Modules, Chari-Venkatesh Modules and Fusion Products",
%   SIGMA 10 (2014), 110, DOI 10.3842/SIGMA.2014.110.
%
% The author's pinned TeX file declares no theorem environments of its own: they are
% supplied by the SIGMA publisher class sigma.cls, which is not installed here.  The
% block below is that class's theorem-environment block (same environment names and
% styles), copied mechanically so that the retained author text compiles unchanged.
% No mathematical macro or statement is introduced by this file.
\newtheorem{theorem}{Theorem}
\newtheorem*{theorem*}{Theorem}
\newtheorem{lemma}{Lemma}
\newtheorem*{lemma*}{Lemma}
\newtheorem{corollary}{Corollary}
\newtheorem*{corollary*}{Corollary}
\newtheorem{proposition}{Proposition}
\newtheorem*{proposition*}{Proposition}
\newtheorem{conjecture}{Conjecture}
\newtheorem*{conjecture*}{Conjecture}
{\theoremstyle{definition} \newtheorem{definition}{Definition}
\newtheorem*{definition*}{Definition}
\newtheorem{example}{Example}
\newtheorem*{example*}{Example}
\newtheorem{remark}{Remark}
\newtheorem*{remark*}{Remark}
\newtheorem{note}{Note}
\newtheorem*{note*}{Note}
}
\usepackage{cleveref}
\crefname{originalreference}{original reference}{original references}
\Crefname{originalreference}{Original reference}{Original references}
\makeatletter
\newlabel{garland}{{garland}{source}{}{sourcepointer-1}{}}
\newlabel{garland@cref}{{[originalreference][0][]garland}{[1][0][]source}{}{}{}}
\makeatother
\title{Three relations holding in the Demazure module attached to a multiple of the highest root}
\author{Bhimarthi Ravinder}
\date{doi:10.3842/SIGMA.2014.110; CC BY 4.0}
\begin{document}
\maketitle
\noindent\textit{Editorial note.} All delivered mathematical text is the author's own, copied from the pinned article source sources/sigma-2014-110-demazure-fusion/pinned-original-source.tex; no statement or proof step was rewritten, reordered or removed.\par\medskip
\noindent\textit{Editorial note.} Preamble provenance: the author's tex declares no theorem environments of its own (they come from the SIGMA publisher class); the staged sources/sigma-2014-110-demazure-fusion/retained-preamble.tex reproduces that class's theorem-environment block verbatim and is used as the preamble, with the wrapper class and base packages supplied by the assembler.\par\medskip
\noindent\textit{Editorial note.} Bibliography provenance: neither the counted statement nor its proof contains a citation (this lane rescans the whole span before writing the card), so the card carries no bibliography and no entry was reconstructed.\par\medskip
\noindent\textit{Editorial note.} Selection provenance: this card was opened under the file-based control contract (CONTROL.md), with the mechanical pre-card gate run on its statement and proof interval before the card was written.\par\medskip
\noindent\textit{Allowed standard prerequisites (editorial).}\begin{itemize}
\item The local Weyl module W(lambda) and its quotient the Demazure module D(1, lambda), with the defining relations (w2) and (w2') remembered as the annihilation of the elements x\textasciicircum{}-\_alpha t\textasciicircum{}m acting on the highest weight vector. The author's sections 2 and 3, in which the modules and the relations are set up before the counted statement. Fixes the module D(1, (k+1)theta) in which the three relations of the statement are asserted.
\item The copy of sl\_2 spanned by x\textasciicircum{}+\_alpha, x\textasciicircum{}-\_alpha and alpha\textasciicircum{}dual inside the current algebra, and the standard sl\_2 weight argument. Used inside the counted proof at source lines 469-476. Produces part (1) from the vanishing of the two vectors displayed in the proof.
\item The commutation relation for two negative root vectors acting on the highest weight vector, labelled garland in the article. Invoked inside the counted proof at source line 478 (a navigation-pointer stub, since the relation is a separate item of the same article). Supplies the sum formula which, after applying x\textasciicircum{}+\_\{theta-alpha\}, yields part (2).
\end{itemize}\medskip
\ifdefined\hypertarget\hypertarget{sourcepointer-1}{}\fi
\noindent\textit{Source reference.} Original reference garland: sources/sigma-2014-110-demazure-fusion/pinned-original-source.tex, source lines 315--315. Navigation pointer only: the counted text refers to the item garland of this same article, which is not part of this card. The pointer carries the item's own author label as its number; its content is not reproduced here.\par\medskip
\begin{lemma}
\label{crucial}
Let $k\geq 1$ and $0\leq i \leq k$. The following relations hold in the module $D(1,(k+1)\theta)$:
\begin{enumerate}\itemsep=0pt
\item[$1)$] $(x^{-}_{\alpha} \otimes 1)^{kd_{\alpha}+1}\big(x^{-}_{\theta} \otimes t^{2k+1-i}\big) \overline{w}_{(k+1)\theta}=0$,
$\forall\, \alpha \in R^{+}$, $(\theta|\alpha)=1$;
\item[$2)$] $\big(x^{-}_{\alpha} \otimes t^{k}\big) (x^{-}_{\theta} \otimes t^{2k+1-i}) \overline{w}_{(k+1)\theta} \in
\big\langle \big(x^{-}_{\theta} \otimes t^{2k+2-i}\big) \overline{w}_{(k+1)\theta} \big\rangle$,
$\forall\, \alpha \in R^{+}$, $(\theta|\alpha)=1$;
\item[$3)$] $\big(x^{-}_{\theta} \otimes t^{2k-i}\big) \big(x^{-}_{\theta} \otimes t^{2k+1-i}\big) \overline{w}_{(k+1)\theta}
\in
\big\langle \big(x^{-}_{\theta} \otimes t^{2k+2-i}\big) \overline{w}_{(k+1)\theta} \big\rangle$.
\end{enumerate}
\end{lemma}

\begin{lemma}
\label{g}
Let $k\in\mathbb{N}$.
The following relations hold in the local Weyl module $W((k+1)\theta)$:
\begin{enumerate}\itemsep=0pt
\item[$1)$] $(x^{-}_{\theta}\otimes 1)^{2k+1}\big(x^{-}_{\theta}\otimes t^{2k+1-i}\big)w_{(k+1)\theta}=0,
\qquad
\forall\,
0\leq i \leq k$;
\item[$2)$] $(x^{-}_{\theta}\otimes t^m)(x^{-}_{\theta} \otimes t^{m+1})  w_{(k+1)\theta} \in  \big\langle \big(x^{-}_{\theta}
\otimes t^{m+2}\big) w_{(k+1)\theta} \rangle,
\qquad
\forall\,
m\geq k$.
\end{enumerate}
\end{lemma}

\begin{proof}
Let $ \alpha \in R^{+} \textup{with} (\theta|\alpha)=1$. This implies that $\theta-\alpha$ is also a~root of
$\mathfrak{g}$ and $(\theta|\theta-\alpha)=1$. We now prove part (1).
Observe that, $(x^{-}_{\theta} \otimes t^{2k+1-i}) \overline{w}_{(k+1)\theta}$ is an element of weight $k\theta$.
Further $(x^{+}_{\alpha} \otimes 1) (x^{-}_{\theta} \otimes t^{2k+1-i}) \overline{w}_{(k+1)\theta}=0$, since
$(x^{+}_{\alpha} \otimes 1)\overline{w}_{(k+1)\theta}=0$ and $(x^{-}_{\theta-\alpha} \otimes t^{2k+1-i}) \overline{w}_{(k+1)\theta}=0$,
for all $0\leq i \leq k$.
Considering the copy of $\mathfrak{sl}_2$ spanned by $x^{+}_{\alpha} \otimes 1$, $x^{-}_{\alpha} \otimes 1$,
$\alpha^{\vee}\otimes 1$, we obtain part (1) by standard $\mathfrak{sl}_2$ arguments.
We now prove part (2).
Putting $r=2$, $s=(3k+1-i)$ and $\alpha=\theta$ in Lemma~\ref{garland}, we get
\begin{gather}
\big(x^{-}_{\theta} \otimes t^{k}\big) \big(x^{-}_{\theta} \otimes t^{2k+1-i}\big) \overline{w}_{(k+1)\theta} +
\sum\limits_{\substack{k+1\leq p\leq q \leq 2k-i\\p+q= 3k+1-i}}
\frac{1}{(2\delta_{p, q})!} \big(x^{-}_{\theta} \otimes t^{p}\big) \big(x^{-}_{\theta} \otimes t^{q}\big)
\overline{w}_{(k+1)\theta}
\nonumber\\
\qquad{}
\in
\big\langle\big(x^{-}_{\theta}\otimes t^{2k+2-i}\big) \overline{w}_{(k+1)\theta} \big\rangle, \label{e}
\end{gather}
since $(x^{-}_{\theta}\otimes 1)^{3k+3-i}\overline{w}_{(k+1)\theta}=0$, $\forall\, 0\leq i \leq k$.
Now we act on both sides of~\eqref{e} by $x^{+}_{\theta-\alpha}$ and use the relation $(x^{-}_{\alpha} \otimes t^r) \overline{w}_{(k+1)\theta}=0$,
for all $r\geq(k+1)$, which gives part~(2).
Part~(3) is immediate from the part~(2) of Lemma~\ref{g}.
\end{proof}

\end{document}
